\section{Extended Ablation Analysis}
\label{sec:ablation_analysis_appendix}

All variants train on mixed HumanoidArena tasks in SONIC X7 and evaluate on DoubleDesk, Football, and P\&Pbox. Figure~\ref{fig:training_loss_comparison} reports mean success across these tasks. Each component ablation is compared separately with the 51.67\% baseline without VLA pretraining; the ablations are not cumulative.

\textbf{Motion prior and backbone freezing.}
Removing the pretrained T2M backbone reduces success to 18.33\%, supporting its contribution to whole-body control. Separately, unfreezing the backbone yields 31.67\%, suggesting that preserving pretrained motion representations is beneficial under this adaptation protocol. These comparisons support retaining a frozen motion generator while learning visual and task conditioning through dedicated trainable branches.

\textbf{Prediction objective and gradient decoupling.}
Replacing MSE with Kimodo's original motion objective yields 35.00\% success, compared with 51.67\% for the baseline. This favors direct supervision of policy prediction targets while retaining the pretrained motion prior. Removing root--body gradient detachment instead yields 45.00\%. Detachment preserves root-to-body conditioning but prevents the body loss from updating the root ControlNet, allowing each stage to learn through its own objective while maintaining coordinated generation. The separate comparisons support both the prediction objective and controlled cross-stage gradient flow.

\textbf{VLA pretraining.}
VLA pretraining improves success from 51.67\% to 56.11\%, with generally lower and less variable training loss. Both variants retain the pretrained T2M backbone, so this improvement indicates a complementary benefit from learning visual and task conditioning. Together, the ablations support the evaluated design choices, but do not measure interactions among simultaneous component changes or establish that their individual gains would combine additively.
